\chapter{完備順序体の基本性質}\label{ch:ofield}

\section{公理 (R1)--(R16)}

第\ref{ch:reals}章で定義した記号 $\R, 0, 1, +, \cdot, -, {}^{-1}$ (添字 $\R$ は省略する)と第\ref{ch:zfc}章の $<$ について, 次はすべて $\ZFC$ の定理である(定理\ref{thm:complete-ordered-field}). 本章以降の証明木では, これらを付録の\textbf{基礎定理}として葉に用いる. 各式は全称閉包を表す.

\begin{center}\small
\begin{adjustbox}{max width=\linewidth}
\begin{tabular}{ll}
  \toprule
  (R1) & $0 \in \R$, \; $1 \in \R$, \; $x \in \R \to y \in \R \to x + y \in \R$ ($\cdot$ も同様), \; $x \in \R \to -x \in \R$, \; $x \in \R \to x \neq 0 \to x^{-1} \in \R$ \\
  (R2) & $x, y, z \in \R \to (x + y) + z = x + (y + z)$ \\
  (R3) & $x, y \in \R \to x + y = y + x$ \\
  (R4) & $x \in \R \to x + 0 = x$ \\
  (R5) & $x \in \R \to x + (-x) = 0$ \\
  (R6) & $x, y, z \in \R \to (x y) z = x (y z)$ \\
  (R7) & $x, y \in \R \to x y = y x$ \\
  (R8) & $x \in \R \to x \cdot 1 = x$, \; $0 \neq 1$ \\
  (R9) & $x \in \R \to x \neq 0 \to x \cdot x^{-1} = 1$ \\
  (R10) & $x, y, z \in \R \to x (y + z) = x y + x z$ \\
  (R11) & $\neg (x < x)$ \hfill (集合一般の定理 [O1] として完全展開) \\
  (R12) & $x < y \to y < z \to x < z$ \hfill (集合一般の定理 [O2] として完全展開) \\
  (R13) & $x, y \in \R \to x < y \lor x = y \lor y < x$ \\
  (R14) & $x, y, z \in \R \to x < y \to x + z < y + z$ \\
  (R15) & $x, y \in \R \to 0 < x \to 0 < y \to 0 < x y$ \\
  (R16) & $X \subseteq \R \to X \neq \emptyset \to \exists b\,\UB(b, X) \to \exists s\,\IsSup(s, X)$ \\
  \bottomrule
\end{tabular}
\end{adjustbox}
\end{center}
(ここで $x, y \in \R \to \cdots$ は $x \in \R \to y \in \R \to \cdots$ の略. 正確な論理式は付録の一覧にある.) 減法と $2$ は定義公理 $x - y = x + (-y)$, $2 = 1 + 1$ で導入し, $x/2$ は $x \cdot 2^{-1}$ のことである.

本章の補題はすべて完全展開した証明木をもつ. 証明木の葉に現れる「$t \in \R$」の証明(例えば $x + (-y) \in \R$)は (R1) の例化と MP で作られ, 木の中にそのまま現れる.

\section{順序の一般的性質}

$<$ は $\subsetneq$ として定義されているので, 次の [O1]--[O7] は実数に限らずすべての集合について成り立つ. 特に $\omega$ の元(自然数)にも実数にも使える.

\begin{lem}\label{lem:order-general}
  \begin{enumerate}[label={[O\arabic*]}]
    \item $\neg\, x < x$.
    \item $x < y \to y < z \to x < z$.
    \item $x < y \to y \leq z \to x < z$.
    \item $x \leq y \to y < z \to x < z$.
    \item $x \leq y \to y \leq z \to x \leq z$.
    \item $x \leq y \to y \leq x \to x = y$.
    \item $x \leq y \to \neg\, y < x$.
    \item[{[O8]}] $x, y \in \R \to \neg\, x < y \to y \leq x$ \quad ((R13) を用いる).
  \end{enumerate}
\end{lem}
\begin{proof}[通常の証明]
  [O1] $x < x$ なら $x \neq x$ となり (A5) に反する. [O2] $x \subseteq y \subseteq z$ より $x \subseteq z$. $x = z$ なら $y \subseteq x$ となり外延性公理から $x = y$ となって $x \neq y$ に反する. [O3]--[O7] は $\leq$ の定義による場合分けと [O1][O2] による. [O8] (R13) で場合分けする.
\end{proof}

\trees{LtIrrefl,LtTrans,LtLeTrans,LeLtTrans,LeTrans,LeAntisym,LeNotLt,NotLtLe}
\section{体の性質}

\begin{lem}\label{lem:field}
  以下で現れる変数はすべて $\R$ の元とする.
  \begin{enumerate}[label={[F\arabic*]}]
    \item $x - y \in \R$.
    \item $2 \in \R$.
    \item (簡約律) $x + z = y + z \to x = y$.
    \item $x + y = 0 \to y = -x$.
    \item $-(-x) = x$.
    \item $-0 = 0$.
    \item $-(x + y) = -x + -y$.
    \item $x \cdot 0 = 0$.
    \item $(-x) y = -(x y)$.
    \item $x (-y) = -(x y)$.
    \item $(-x)(-y) = x y$.
    \item $(x - y) + y = x$.
    \item $x - x = 0$.
    \item $-(y - x) = x - y$.
    \item $(x + y) - (a + b) = (x - a) + (y - b)$.
    \item $a - b = (a - x) + (x - b)$.
    \item $c x - c a = c (x - a)$.
    \item $x y - a b = x (y - b) + b (x - a)$.
  \end{enumerate}
\end{lem}
\begin{proof}[通常の証明]
  [F3] $x = x + 0 = x + (z + (-z)) = (x + z) + (-z) = (y + z) + (-z) = y + (z + (-z)) = y + 0 = y$. [F4] $y = y + 0 = y + (x + (-x)) = (y + x) + (-x) = (x + y) + (-x) = 0 + (-x) = -x$. [F5]--[F7] は [F4] による. [F8] $x \cdot 0 + x \cdot 0 = x(0 + 0) = x \cdot 0 = 0 + x \cdot 0$ と [F3]. [F9] $x y + (-x) y = y x + y(-x) = y(x + (-x)) = y \cdot 0 = 0$ と [F4]. 他も (R2)--(R10) による同様の等式変形である(各ステップは証明木の中で (A6) による書き換えとして現れる).
\end{proof}

\trees{SubR,TwoR,AddCancel,NegUnique,NegNeg,NegZero,NegAdd,Mul0,NegMul,MulNeg,NegNegMul,SubAdd,SubSelf,SubNeg,SubAddAdd,SubSplit,MulSub,MulSubSplit}
\section{順序と演算}

\begin{lem}\label{lem:order}
  以下で現れる変数はすべて $\R$ の元とする.
  \begin{enumerate}[label={[F\arabic*]}, start=19]
    \item $x \leq y \to x + z \leq y + z$.
    \item $x < y \to u \leq v \to x + u < y + v$.
    \item $x \leq y \to u \leq v \to x + u \leq y + v$.
    \item $x < y \to -y < -x$.
    \item $x \leq y \to -y \leq -x$.
    \item $x < y \to 0 < y - x$.
    \item $0 < y - x \to x < y$.
    \item $x < y \to x \neq y$ (集合一般).
    \item $(-1)(-1) = 1$.
    \item $0 < 1$.
    \item $0 < 2$.
    \item $2 \neq 0$.
    \item $2^{-1} \in \R$.
    \item $0 < z \to x < y \to x z < y z$.
    \item $0 \leq z \to x \leq y \to x z \leq y z$.
    \item $0 < x \to 0 < x^{-1}$.
    \item ($\varepsilon/2$) $0 < \varepsilon \to 0 < \varepsilon/2$.
    \item ($\varepsilon/2$) $\varepsilon/2 + \varepsilon/2 = \varepsilon$.
  \end{enumerate}
\end{lem}
\begin{proof}[通常の証明]
  [F22] $x < y$ の両辺に $(-x) + (-y)$ を加え ((R14)), 両辺を整理する. [F28] (R13) で $0, 1$ を比較する. $0 = 1$ は (R8) に反する. $1 < 0$ なら [F22] で $0 < -1$, (R15) で $0 < (-1)(-1) = 1$ となり [O2][O1] に反する. [F34] $x^{-1} = 0$ なら $1 = x x^{-1} = 0$. $x^{-1} < 0$ なら $0 < x(-x^{-1}) = -1$ となり $-1 < 0$ と矛盾. [F36] $\varepsilon/2 + \varepsilon/2 = (\varepsilon/2)(1 + 1) = \varepsilon(2^{-1} \cdot 2) = \varepsilon$.
\end{proof}

\trees{LeAddRight,AddLtLe,AddLeLe,NegLt,NegLe,SubPos,PosSub,LtNe,NegOneSq,ZeroLtOne,TwoPos,TwoNe,InvTwoR,MulLtPos,MulLeNonneg,InvPos,HalfPos,HalfSum}
\section{自然数の埋め込み}

\begin{dfn}
  再帰定理(定理\ref{thm:recursion-omega})により, $\iota : \omega \to \R$ を $\iota(0) = 0$, $\iota(S(n)) = \iota(n) + 1$ で定める. $n \in \omega$ に対して $\iota(n)$ を $\nat{n}$ と書く(テキストの「$na$」は $\nat{n} \cdot a$ である). $1_\omega := S(\emptyset)$ (自然数としての1) と書く.
\end{dfn}

$\iota$ について用いる性質は付録の基礎定理 [N4]--[N6] ($\nat n \in \R$, $\nat{1_\omega} = 1$, $\nat{S(n)} = \nat n + 1$)である. いずれも再帰定理による $\iota$ の定義と, $n$ に関する帰納法から従う.


\medskip\noindent\textbf{帰納法の構成手続き.}\quad
$\N := \omega \setminus \{\emptyset\}$ (定義公理 $n \in \N \leftrightarrow n \in \omega \land n \neq \emptyset$)上の帰納法「$\varphi(1_\omega)$ かつ $\forall n \in \N\,(\varphi(n) \to \varphi(S(n)))$ ならば $\forall n \in \N\,\varphi(n)$」は, 定理\ref{thm:induction}の証明を $\varphi$ ごとにそのまま繰り返す手続きとして用いる(証明生成器 \texttt{prover/p\_nat.js} の \texttt{indOmega}, \texttt{indNat}). すなわち, 帰納法を用いる各定理の証明木の中には, 分出公理 \textsf{Sep} による集合 $\{n \in \omega \mid \varphi(n)\}$ の存在と, それが帰納的であることの証明が毎回展開されて現れる. 単調性「$m \leq n \to a_m \leq a_n$」も同様の手続き \texttt{monoBuild} で示す.

\begin{lem}[自然数と $\iota$]\label{lem:nat}
  \lemlist{NatSub,SNe0,NatSO,NatOne,NatS,NatOmega,LeS,OneLe,LtAddOne,IotaGeOne,IotaPos,IotaMono}
\end{lem}
\begin{proof}[通常の証明]
  [N7][N0]--[N2][N7'][N8] は定義と定理\ref{thm:omega-order}から. [N12] $1_\omega \leq n$, [N14] $1 \leq \nat n$ は $\N$ 上の帰納法, [N15] は [N14] と [F28], [N16] $m \leq n \to \nat m \leq \nat n$ は \texttt{monoBuild} による.
\end{proof}

\section{順序体の補題(追加)}

第II部で繰り返し用いる順序体の計算規則をまとめる. いずれも (R1)--(R16) と前節までの補題から通常の計算で示される. 第II部の各章にも, その章でのみ用いる補題 [E28]--[E51] がある.

\lemlist{LeAddNonneg,LeAddNonnegL,LtAddPos,LtAddCancel,AddPosPos,AddNonnegNonneg,AddLtLt,MulLtPosL,MulLeNonnegL,PosNonneg,MulPosPos,MulLeLe,InvUnique,InvNe0,MulNe0,InvMul,Abs1,AbsInv,InvLt,InvLe,SubInv,InvMulL,DivMul,LtDiv,MulLtDiv,AbsLowerHalf,RevTriangle}
